\textbf{Proof.}
Write \(L=L_n=\log(en)\), \(m=m_2\), and \(f=(\Lambda+\Gamma)/2\).  The proof has four steps.

\emph{Step 1: support, smooth factors, and the pilot event.}
For \(u\in\mathcal C_\epsilon\setminus\{0\}\),
\[
 s_0(u)\ge\epsilon s(u),\qquad s_1(u)\ge\epsilon s(u).
\]
Consequently all divisions defining \(\alpha\) and \(\beta\) are valid, \(0\le\alpha(u),\beta(u)\le s(u)\), and
\[
 f(u)=\frac12\{\alpha(u)+\beta(u)
  +(\sqrt{\alpha(u)}+\sqrt{\beta(u)})
     \lvert\sqrt{\alpha(u)}-\sqrt{\beta(u)}\rvert\}.
\]
All three functions are degree-one homogeneous, including at the cone vertex under the stipulated zero convention.

For a multinomial coordinate, Bernstein's inequality gives, for a numerical \(c>0\),
\[
 P\!\left(
  \left\lvert\widehat q^{(0)}_{ay,x}-q_{ay,x}\right\rvert
  >8\sqrt{q_{ay,x}L/m_0}+16L/m_0
 \right)\le2e^{-cL}.
\]
The constants in \(w_{ay,x}\) dominate the right-hand radius after replacing the unknown coordinate by its empirical version.  A union bound over four coordinates and at most \(d\le\rho_\epsilon nL\) cells, with \(\rho_\epsilon\) reduced if necessary and the finitely many small \(n\) absorbed into \(C_\epsilon\), produces an event \(\mathcal G_0\) with
\[
 P(\mathcal G_0^c)\le C_\epsilon n^{-4},
 \qquad q_x\in B_{x,n}\quad\hbox{for every }x.
\]
The same inequalities applied to arm totals and Bernoulli outcome sums give the pilot sandwich: on \(\mathcal G_0\), membership in \(\widehat{\mathcal R}_n\) implies \(p_x\ge c_\epsilon L/n\) and \(\lvert\Delta_x\rvert\ge c_\epsilon\sqrt{L/(np_x)}\), while every cell violating either condition is sent to \(\widehat{\mathcal D}_n\).  No assertion is needed on null cells; there \(q_x=0\), and zero anchoring makes their contribution vanish.

\emph{Step 2: polynomial approximation on the two pilot regions.}
On a sparse box, rescale the four coordinates by its pilot noise radius.  The fixed Chebyshev reciprocal polynomial applied to \(s(u)u_{11}/s_1(u)+s(u)u_{01}/s_0(u)\), together with the displayed square-root factorization, yields a polynomial of total degree at most \(K_n\).  The zero-constant construction and the coefficient bounds in \(\mathrm{lem:jhw-local-polynomial-toolkit}\), after the overlap inequalities above change the scale by at most \(\epsilon^{-1}\), give
\[
 \lvert P^\star_{x,n}(u)-f(u)\rvert
 \le C_\epsilon\left\{
  \sqrt{\frac{s(u)}{nL}}+\frac1{nL}
 \right\},
 \qquad u\in B_{x,n}.
\]
The rational rounding adds at most \(n^{-3}\); since \(n^{-3}\le(nL)^{-1}\), it is absorbed.  The subtraction at zero proves \(P^\star_{x,n}(0)=0\).

On a heavy near-tie box, the arm denominators are at least \(c_\epsilon p_x\), so the ratio map is smooth in the standardized coordinate directions.  Only the absolute-value cusp remains.  Centering at \(\widehat q_x^{(0)}\), the box has contrast width at most \(C_\epsilon W_{x,n}\).  The degree-\(K_n\) Chebyshev approximation to \(\lvert t\rvert\) has uniform error at most \(C W_{x,n}/K_n\).  Tensorization with the three smooth directions and rational rounding therefore gives
\[
 \sup_{u\in B_{x,n}}
 \lvert P^\star_{x,n}(u)-f(u)\rvert
 \le C_\epsilon W_{x,n}/K_n+C_\epsilon n^{-3}.
\]
Since \(W_{x,n}/K_n\le C_\epsilon\sqrt{p_x/(nL)}+C_\epsilon/(nL)\) on \(\mathcal G_0\), both regions obey the same summable envelope at the true \(q_x\).  Cauchy--Schwarz and \(\sum_xp_x=1\) yield
\[
 \begin{aligned}
 \sum_{x\in\widehat{\mathcal D}_n}
  \lvert P^\star_{x,n}(q_x)-f(q_x)\rvert
 &\le C_\epsilon\left\{
   \sqrt{\frac d{nL}}+\frac d{nL}\right\},\\
 \left(\sum_{x\in\widehat{\mathcal D}_n}
  \lvert P^\star_{x,n}(q_x)-f(q_x)\rvert\right)^2
 &\le C_\epsilon\frac d{nL},
 \end{aligned}
\]
where the last inequality uses \(d\le\rho_\epsilon nL\) and \(\rho_\epsilon\) is chosen sufficiently small.

\emph{Step 3: factorial variance, including cross-cell covariance.}
Condition on \(I_0\).  The polynomial coefficients are then fixed.  For orientation, the fixed-multinomial lift in \(\mathrm{lem:factorial-lift-unbiased}\) obeys \(E(U_{x,r}\mid I_0)=q_x^r\), and expansion into ordered distinct sample indices gives, for \(\lvert r\rvert,\lvert s\rvert\le K_n\),
\[
 E(U_{x,r}U_{x,s}\mid I_0)
 \le e\,q_x^r\left(q_x+\frac{K_n}{m}\mathbf 1_4\right)^s.
\]
For \(x\ne z\), disjoint multinomial categories give exactly
\[
 E(U_{x,r}U_{z,s}\mid I_0)
 =\frac{(m)_{\lvert r\rvert+\lvert s\rvert}}
 {(m)_{\lvert r\rvert}(m)_{\lvert s\rvert}}
 q_x^rq_z^s,
\]
and, when \(2K_n\le m/2\),
\[
 \left\lvert
  \frac{(m)_{\lvert r\rvert+\lvert s\rvert}}
  {(m)_{\lvert r\rvert}(m)_{\lvert s\rvert}}-1
 \right\rvert
 \le\frac{2K_n^2}{m}.
\]
These identities are algebraic consequences of cancelling falling factorials.  Inserting the enlarged-box majorants \(A_{x,n}\) into the last display shows explicitly where the crude fixed-multinomial cross bound can retain the factor \(K_n^2/m\).  The corrected construction removes that loss by fixed-fraction Poisson thinning.

Indeed, in the uncapped comparison experiment let \(M\sim\operatorname{Poisson}(\lambda_2)\), \(\lambda_2=m/2\), independently of an infinite continuation of the split-two stream.  Poisson splitting makes all category counts \(\overline Z_{ay,x}\) mutually independent with means \(\lambda_2q_{ay,x}\).  Hence
\[
 E\left[\frac{(\overline Z_{ay,x})_k}{\lambda_2^k}\right]
 =q_{ay,x}^k,
 \qquad
 \operatorname{Cov}\left(
  \overline U_{x,r},\overline U_{z,s}\right)=0
 \quad(x\ne z).
\]
The zero constant on sparse boxes removes the empty-cell contribution, while the centered stripe basis supplies one pilot noise radius for every nonconstant term.  Applying the cited coefficient estimates at the enlarged point \(q_x+K_n\mathbf1_4/\lambda_2\), and using \(K_n\le\epsilon^2L/2^{16}\), gives the within-cell sum
\[
 \sum_{x\in\widehat{\mathcal D}_n} E\left[
  \left\{
   \sum_r c_{r,x}(\overline U_{x,r}-q_x^r)
  \right\}^2\middle|I_0\right]
 \le C_\epsilon\left\{\frac1n+\frac d{nL}\right\}.
\]
Independence makes every cross-cell covariance exactly zero before capping.

The Poisson Chernoff bound at twice the mean gives
\[
 P(M>m)\le\exp(-m/6).
\]
Both the randomized preliminary estimate and the target are projected to \([0,1]\), so replacing the uncapped estimate on \(\{M>m\}\) changes squared risk by at most \(C\exp(-m/6)\).  Finally, for the conditional-expectation estimator in the corrected definition, conditional Jensen gives
\[
 \left\{E(\widetilde V\mid O_1,\ldots,O_n)-V^\star(P)\right\}^2
 \le E\left[\{\widetilde V-V^\star(P)\}^2\mid O_1,\ldots,O_n\right].
\]
Thus the nonrandomized fixed-sample estimator inherits the uncapped within-cell bound, zero cross-cell covariance, and exponentially small cap error.  This accounts for the full covariance after derandomization without imposing an additional sampling assumption.

\emph{Step 4: regular cells and pilot mistakes.}
On \(\widehat{\mathcal R}_n\cap\mathcal G_0\), the pilot separation is larger than the split-one ratio fluctuation.  A further Bernstein bound shows that the maximizing arm is wrong with probability at most \(C(en)^{-8}\) per cell.  On the complementary event the selected contribution is a smooth ratio functional with arm denominators at least \(c_\epsilon p_x\).  Conditional variance expansion and \(\sum_xp_x=1\) give
\[
 E\left[
  \left\{\sum_{x\in\widehat{\mathcal R}_n}
   (\widehat v_x^{\mathrm{rat}}-f(q_x))\right\}^2
  \middle| I_0\right]
 \le C_\epsilon/n+C_\epsilon n^{-4}.
\]
The union of wrong-arm, box-failure, and pilot-sandwich errors has probability at most \(C_\epsilon n^{-4}\).  Since both the target and the final projected estimator lie in \([0,1]\), its total squared-loss contribution is at most that probability.  Combining the four steps proves every assertion.